\begin{table}[htbp]
\centering
\caption{\textbf{Target recovery under one evaluation framework.} Transcription factors with significant target enrichment (FDR $< 0.05$).}
\label{tab:ceiling}
\smallskip
\small
\begin{tabular}{lccc}
\toprule
Method & TRRUST & DoRothEA A+B+C & DoRothEA (all) \\
\midrule
Differential expression of the knockdown & 1/9 (11.1\%) & 3/13 (23.1\%) & 8/18 (44.4\%) \\
SAE feature response & 0/9 (0.0\%) & 0/13 (0.0\%) & 0/18 (0.0\%) \\
GENIE3 & 1/9 (11.1\%) & 1/13 (7.7\%) & 0/18 (0.0\%) \\
GRNBoost2 & 1/9 (11.1\%) & 1/13 (7.7\%) & 1/18 (5.6\%) \\
Co-expression (Pearson) & 2/9 (22.2\%) & 2/13 (15.4\%) & 2/18 (11.1\%) \\
Co-expression (Spearman) & 2/9 (22.2\%) & 2/13 (15.4\%) & 1/18 (5.6\%) \\
Geneformer gene embeddings & 1/9 (11.1\%) & 1/13 (7.7\%) & 0/18 (0.0\%) \\
Random gene sets & 0/9 (0.0\%) & 0/13 (0.0\%) & 0/18 (0.0\%) \\
\bottomrule
\end{tabular}
\\[2pt]
\footnotesize Each method emits a ranked list of 100 putative targets per transcription factor; enrichment for that factor's curated targets is tested by hypergeometric test with Benjamini--Hochberg correction across the panel, with the universe set to the genes the assay measures. Differential expression of the knockdown is perturbation-aware and bounds what these data support.
\end{table}
